Recorded sales are a censored measurement of demand: stockout days record zeros and capacity days record caps, and forecasting models fit to raw sales learn the censoring as if it were demand, feeding under-orders that cause the next stockout. This practitioner's guide presents three probabilistic mechanisms for forecasting demand rather than sales, selected by what the data records about each stockout. A binary on-shelf flag admits \emph{gating}: an intermittent-demand smoother in the Croston and Teunter-Syntetos-Babai family freezes its occurrence updates during off-shelf periods. A per-observation censoring point admits a \emph{censored likelihood}: capped observations enter a Tobit-style model as lower bounds, which halves peak-day errors in simulation exactly where a plain likelihood saturates at the cap. A fractional, noisily labeled availability admits a \emph{multiplicative saturating factor with a learned floor}, demonstrated on the FreshRetailNet-50K dataset, where the full $50{,}000$-series hierarchical model fits in under six minutes of stochastic variational inference on a GPU. All three share one structure: the corruption is a switchable model component, and the demand forecast is the model with the component switched off. We pair the mechanisms with an evaluation discipline for censored settings, showing concretely that aggregate accuracy against recorded sales rewards reproducing the corruption, and we state the limitations of each mechanism explicitly, including exogenous-availability assumptions, Gaussian misspecification, and the absence of cross-mechanism comparisons. All case studies are reproducible from the open-source \pkg{numpyro\_forecast} package.
